% =====================================================================
\section{Медленная защита: что добавилось и что она сделала}\label{sec:slow}
\figp{fig_slow.png}{\label{fig:slow}Четыре результата версии~4. \textbf{а}~--- две ступени защиты: сигнал индукции быстрый, белок медленный, и~половина уровня набирается только к~третьим суткам. \textbf{б}~--- закон длительности эпизода: он оказался прямой, а~не~насыщающейся кривой, и~проверочная точка вне подгонки это подтвердила. \textbf{в}~--- ритм ломается не~плавно, а~удвоениями длины цикла. \textbf{г}~--- число клеток входит в~модель как эффективный разброс стойкости.}

Закон сохранения разрывов (раздел~\ref{sec:law}) оставил ровно два выхода, и~второй из~них~--- переменная защита: если клетка после нагрузки усиливает ремонт, нижняя граница окна подтягивается вверх сама, и~возврат не~требует большого выброса шума. Этот выход мы и~достраивали.

\textbf{Что именно добавлено.} Фермент ремонта пероксиредоксина (сульфиредоксин) наращивается не~мгновенно: сначала накапливается сигнал индукции, потом по~нему синтезируется белок. Это две ступени первого порядка подряд:
\[\dot m=\frac{1}{\tau_m}\Big(\beta\frac{u}{K_B+u}-m\Big),\qquad \dot s=\frac{1}{\tau_s}\,(m-s),\]
где $u$~--- перекись, $m$~--- сигнал, $s$~--- наработанный белок, который прибавляется к~постоянной защите. Стационар тот~же, что был раньше: окно, его границы и~весь закон сохранения разрывов не~меняются. Меняется только ход во~времени~--- и~именно он решает всё.

\textbf{Откуда постоянные.} Время первой ступени взято по~мРНК: 8,1--9,6~ч в~трёх независимых обработках.\tS\ Время второй~--- по~белку из~сырых данных той~же лаборатории: пик на~18-м часу, а~постоянная спада в~опубликованных работах не~измерена ни~разу, и~по~разным способам обработки выходит от~24 до~88~ч.\tS\ Рабочее значение 63~ч~--- середина этого разброса, и~это наше допущение, а~не~измерение.\tD\ Сила индукции разыграна по~наборам: измеренная кратность наработки белка идёт $\times$2,1 / 2,7 / 4,2 / 7,0 при~четырёх дозах окислителя.\tS

\textbf{Что из~этого вышло~--- и~это главный отрицательный результат версии.} Длительность эпизода оказалась \emph{прямой} функцией памяти защиты (рис.~\ref{fig:slow}б):
\[\text{длительность}\;=\;1{,}28\;+\;0{,}97\,\tau_{\rm б}\quad\text{(суток)},\]
без всякого насыщения. Форма выбиралась не~на~глаз: три кандидата подгонялись по~точкам с~памятью ниже 100~ч, а~проверялось предсказание в~точке 150~ч, в~подгонку не~входившей. Прямая выиграла у~обеих насыщающихся и~по~невязке подгонки, и~по~этой проверке.

Отсюда следствие, которое закрывает целую дорогу. Чтобы эпизод уложился в~клинические двое суток, нужна память защиты около 12~часов~--- то~есть \emph{меньше, чем у~первой ступени}, и~в~несколько раз меньше любого из~измеренных значений (24--88~ч). Пассивная защита не~промахивается мимо цели на~волосок: она исключена в~2--10~раз. Усиление ремонта само по~себе эпизод не~укорачивает, а~удлиняет~--- потому что чем дольше держится наработанный белок, тем дольше клетка сидит в~промежуточном состоянии, вместо того чтобы быстро свалиться обратно.

\begin{ineq}
Две ступени первого порядка подряд дают отклик на~ступень нагрузки
\[s(t)=1-\frac{\tau_s e^{-t/\tau_s}-\tau_m e^{-t/\tau_m}}{\tau_s-\tau_m},\]
то~есть всегда с~перегибом: мгновенного отклика нет ни~при каких постоянных. При~$\tau_m=12$~ч и~$\tau_s=63$~ч половина уровня набирается к~57~часам (рис.~\ref{fig:slow}а). Это и~есть причина, по~которой ускорять надо не~первую ступень, а~\emph{усиление}: производная отклика по~$\tau_m$ в~рабочей точке мала (на~одной клетке ускорение дало $-1\dots-5$~\%), а~по~амплитуде~--- велика ($-52$~\%). Предсказание было записано до~счёта и~подтвердилось.
\end{ineq}

\begin{note}
Из~того, что защита в~модели эпизод удлиняет, \emph{не} следует, что сульфиредоксин вреден. Следует одно: защита, которая только чинит повреждённое, картину болезни не~объясняет. Объяснить её может лишь защита, которая срабатывает \emph{раньше и~сильнее}~--- то~есть примирование (раздел~\ref{sec:prime}).
\end{note}

% =====================================================================
\section{Примирование: единственная дорога, которая довела до~цели}\label{sec:prime}
Примирование~--- это память о~прошлом воздействии: клетка, уже встречавшая окислитель, отвечает на~следующее воздействие сильнее и~быстрее. В~модель оно вставлено одной строкой: тем~же сигналом, что движет индукцию, набирается след $p$, и~сила индукции умножается на~$(1+A\,p)$. Потолок $A$ взят по~измерению~--- втрое.

\facts
\begin{itemize}
\item При~повторном воздействии через 24~часа отклик регулятора ниже, а~наработка самого фермента \emph{выше}~--- до~трёх раз, и~время до~полумаксимума падает примерно с~8 до~4~часов. Та~же лаборатория и~та~же система, что дала нам сырые данные по~кинетике.\tS
\item Эффект зависит от~главного регулятора окислительного ответа и~не~зависит от~его вспомогательного партнёра: то~есть это именно усиление сигнала, а~не~побочный эффект.\tS
\item Независимо, в~другой ткани: накопление защитного фермента при~повторных воздействиях прослежено на~четырёх сутках и~избирательно по~генам.\tP
\end{itemize}

\textbf{Что получилось.} Впервые за~проект длительность эпизода вошла в~клиническую цель: медиана 1,46~суток при~верхней границе 1,665 (односторонняя оценка по~порядковой статистике), доля коротких эпизодов 0,363, парный критерий против контроля без примирования $p=5{,}2\cdot10^{-6}$. Все четыре части правила, записанного до~счёта, выполнены~--- и~повторно выполнены на~выборке в~2000~наборов, вшестеро большей.

Важно, чем это \emph{не} является. Примирование не~подгонялось: его амплитуда взята потолком из~измерения, его постоянная времени перебиралась, а~ускорение первой ступени~--- вторая половина измеренного эффекта~--- было проверено отдельно и~не~дало почти ничего, как и~предсказывалось до~счёта. Дорога осталась одна, и~она прошла.

% =====================================================================
\section{Предвестник: кто его делает}\label{sec:prodrome}
За~несколько суток до~срыва бодрствование удлиняется~--- в~модели это видно как подъём кривой «длина бодрствования против срока до~срыва», нормированной на~собственный здоровый уровень каждого набора. Вопрос, кто его делает, решался парным сравнением: одни и~те~же наборы, одно зерно, разница только в~одном звене.

\begin{itemize}
\item \textbf{Медленный шум нагрузки~--- не~он.} При~его включении высота предвестника не~растёт (парный знаковый критерий: подъём у~131 набора из~выборки, $p=0{,}97$). Шум вводился под~обоснование «в~модели нет ни~одной величины, живущей на~масштабе дней»; это перестало быть правдой, как только появилась двухступенчатая защита. Шум выведен из~рабочего набора.
\item \textbf{Память защиты~--- частично он.} При~её отнятии высота падает у~107~наборов из~проверенных, $p\le0{,}05$. Но~уходит только \emph{пятая часть} превышения.
\end{itemize}

Значит есть третья медленная величина. Поиск сузился до~одного кандидата простым перебором всех времён модели: возврат перекиси и~релаксация нагрузки дают 0,25--0,5 и~0,03--0,05~суток; возврат $x$~--- порядка 0,01~суток (и~это подтверждено прямым измерением: у~73~\% наборов $x$ возвращается за~сутки); первая ступень защиты~--- 0,5~суток; вторая~--- 2,6~суток и~уже исключена как единственная причина. \emph{Других времён в~уравнениях нет.}

Отсюда гипотеза, записанная как правило Ф21: третья медленная величина~--- не~постоянная времени, а~\textbf{накопление}. Население из~$J$ клеток заполняется повреждёнными по~одной, и~характерное время заполнения задаётся не~константой, а~\emph{редкостью} перехода, то~есть расстоянием нагрузки до~барьера в~сигмах шума. У~такой величины нет своего~$\tau$~--- поэтому перебор постоянных её найти и~не~мог.

% =====================================================================
\section{Из~чего состоит эпизод}\label{sec:tail}
Все прочие правила меряют длительность \emph{между} наборами. Но~для~выбора механизма важнее, из~чего длительность складывается \emph{внутри} набора. Эпизод делится на~две части: время, пока нагрузка держится выше верха окна, и~\textbf{хвост}~--- время от~последнего такого момента до~конца эпизода, когда нагрузка уже нормальная, а~клетки ещё повреждены.

Это различие решает, за~что вообще может зацепиться защита. Если хвост велик, лимитирует ремонт, и~ускорять надо его. Если хвоста почти нет, эпизод длится ровно столько, сколько держится высокая нагрузка, и~ремонт не~лимитирует вовсе.

На~прежней версии модели (ещё без двухступенчатой защиты) хвост занимал четверть эпизода, а~связь длительности со~временем высокой нагрузки была 0,94. Предварительная проверка на~\emph{нынешней} модели даёт картину резче: связь 0,87, хвост~--- около 7~\% эпизода.\tD\ Если это подтвердится на~полной выборке, вывод будет жёстким: \textbf{примирование укорачивает эпизод не~тем, что чинит быстрее, а~тем, что поднимает верх окна} и~тем сокращает само время высокой нагрузки. Это проверяемое предсказание, и~оно записано до~счёта.

\begin{note}
Прежнее измерение хвоста сделано на~модели без двухступенчатой защиты, и~переносить его на~нынешнюю нельзя: именно защита двигает границы окна. Поэтому правило переписано, а~регистратор эпизодов \emph{сверяется побитово} с~основным счётом на~тех~же строках и~том~же зерне: если копия шага разойдётся хоть в~одном эпизоде, правило объявляет себя непроверяемым и~чисел не~печатает.
\end{note}

% =====================================================================
\section{Сколько клеток в~населении~--- и~почему это оказался параметр}\label{sec:J}
Число клеток $J$ до~недавнего времени выбиралось по~удобству счёта: 32 в~одних прогонах, 64 в~других. Предполагалось, что это просто разрешение. Оказалось, что это параметр, и~входил он тремя разными путями.

\begin{enumerate}
\item \textbf{Хвост распределения стойкости.} Множители брались равновероятными квантилями $p_j=(j+1/2)/J$, поэтому крайний квантиль уползает с~ростом $J$: чем больше клеток, тем уязвимее самая уязвимая. \emph{Закрыто}: сетка вероятностей закреплена по~краям, и~при~$J=32$ она совпадает с~прежней до~последнего знака, так что старые прогоны остаются в~силе.
\item \textbf{Квантование порога эпизода.} Эпизод определялся как «доля повреждённых выше $1/2$», а~доля принимает только значения $j/J$, поэтому настоящий порог равен $0{,}625$ при~$J=8$ и~$0{,}516$ при~$J=64$. \emph{Обойдено}: вердикты выносятся по~непрерывной переменной, у~которой порог не~квантован.
\item \textbf{Реализованный разброс.} А~вот это нашлось только на~выборке в~2000~наборов, когда шум измерений упал ниже эффекта. Края сетки закреплены и~действительно почти не~ползут (наименьший множитель 0,8002 при~$J=8$ и~0,8029 при~$J=512$), но~\emph{дисперсия} конечной сетки меньше заданной и~подходит к~ней только в~пределе: $\mathrm{std}(\ln v)=0{,}1226\,/\,0{,}1060\,/\,0{,}0980\,/\,0{,}0943$ при~$J=8/16/32/64$ и~$0{,}0910$ в~пределе (рис.~\ref{fig:slow}г). От~$J=16$ к~$J=64$ это 11~\% по~эффективной сигме разброса, а~модель к~ней крайне чувствительна.
\end{enumerate}

Третий путь объясняет найденную зависимость количественно и~указывает поправку: растянуть сетку так, чтобы реализованный разброс равнялся заданному при~любом~$J$. Решающая проверка записана до~счёта: блок при~$J=32$ с~подменённой сигмой обязан воспроизвести блок при~$J=16$ с~обычной. Если воспроизведёт~--- путь объяснён целиком.

\begin{ineq}
Общий урок, не~про~нашу модель. Дискретизация непрерывного распределения $J$~точками может \emph{точно} сохранять среднее (нормировкой) и~\emph{точно} закреплять края (выбором крайних квантилей) и~при~этом не~сохранять дисперсию: для~равномерной по~вероятности сетки $\mathrm{Var}\to\sigma^2$ лишь как $O(1/J)$ сверху. Если модель чувствительна к~дисперсии разброса~--- а~здесь чувствительность по~логарифму около~2,5,~--- то~число точек сетки становится параметром модели, и~сравнивать прогоны с~разным $J$ нельзя: это сравнение двух разных моделей.
\end{ineq}

% =====================================================================
\section{Как ломается ритм}\label{sec:bifur}
Отдельный вопрос, поставленный независимо от~редокса: что происходит с~чередованием сна и~бодрствования, когда повреждение всё сильнее тормозит сонное ядро. Ожидание было такое: при~приближении к~критическому значению должен появляться хаос~--- то~есть циркадный ритм ломается \emph{при~сохранённом} 24-часовом внешнем задатчике. Для~релаксационного осциллятора под~периодической силой это штатное поведение.

Считается это дёшево: переключатель детерминированный, клеток и~редокса не~нужно вовсе. Результат~--- лестница захвата (рис.~\ref{fig:slow}в):
\begin{itemize}
\item до~$kD\approx4{,}8$~мВ цикл равен суткам~--- захват 1:1;
\item при~$kD\approx4{,}8$ цикл скачком удваивается до~48~ч: \emph{ночь пропускается};
\item при~$kD\approx5{,}2$ ещё одно удвоение, до~96~ч: пропускаются три ночи;
\item при~$kD\approx5{,}25$ сон исчезает совсем.
\end{itemize}

Это и~есть искомая картина: ритм ломается не~потому, что исчез задатчик, а~потому, что срывается \emph{захват}. Но~рабочая область модели до~неё не~достаёт: при~$k=2$~мВ и~повреждении не~больше единицы потолок равен 2~мВ, то~есть вдвое ниже первой бифуркации. Значит нынешняя модель поломку ритма объяснить не~может, и~это честный отрицательный ответ, а~не~недосмотр.

\begin{note}
Сначала наблюдаемая была выбрана неверно: искалось \emph{множество} длительностей бодрствования, которое перестаёт быть одноэлементным. Но~в~захваченном режиме длительность одна при~любом периоде, а~удвоение видно в~\emph{длине цикла}. Правило, записанное до~счёта, спрашивало не~то; оно сохранено как есть, а~рядом записан двойник с~правильной наблюдаемой. Отношение ширин двух ступеней лестницы вышло около~4~--- близко к~постоянной Фейгенбаума 4,669, но~по~двум ступеням её не~измеряют, и~это пока наблюдение, а~не~измерение.
\end{note}

% =====================================================================
\section{Как мы теперь считаем}\label{sec:method}
Правило фиксируется до~счёта; это было с~самого начала. Версия~4 добавила к~нему вторую половину: \emph{у~каждого критерия должны быть известны его нулевое значение, его шум при~нашем числе наборов и~его мощность}. Иначе вердикт говорит о~выборке, а~не~о~модели.

Формально так. Часть правила считается проверяемой, только если расстояние от~нулевого значения до~порога не~меньше минимального различимого эффекта
\[\mathrm{MDE}=(z_{0{,}95}+z_{0{,}80})\,\mathrm{SE}=2{,}4865\,\mathrm{SE}.\]
Иначе вердикт~--- «непроверяемо», и~\textbf{симметрично}: и~там, где правило провалилось~бы, и~там, где прошло~бы. Рядом печатается, сколько наборов нужно, чтобы порог стал различим. Для~утверждений вида «наблюдение согласуется с~предсказанием» допуск взят в~3~SE: он обязан быть шире MDE, иначе такое утверждение неопровержимо в~принципе.

Это не~формальность: четыре правила оказались непроверяемыми при~любом числе наборов, потому что мерились по~восьми ячейкам, где ошибка ранговой связи равна $1/\sqrt{7}=0{,}38$, а~порог ставился $-0{,}5$. Лечится это не~выборкой, а~дроблением ячеек.

\textbf{Четыре статистические ловушки, найденные по~дороге}, каждая стоила отдельного прогона:
\begin{enumerate}
\item \emph{Медиана по~наборам от~доли редких событий} всегда 0 или~1: она обнуляется, едва события есть меньше чем у~половины наборов. Три правила падали на~этом, а~не~на~модели.
\item \emph{Доля коротких промежутков} тождественно равна $1+\mathrm{CV}^2$ по~частотам наборов и~физики не~видит вовсе. Замена: $\mathrm{CV}^2$ промежутков \emph{внутри} набора, где при~глухом периоде $\mathrm{CV}^2=(1-t_r/g)^2$, а~нуль равен единице ровно.
\item \emph{Направление знакового критерия}: одна и~та~же функция, применённая к~«меньше~--- лучше» и~«больше~--- лучше», дала $p=0{,}99$ там, где правильное направление давало 0,038. Теперь направление~--- обязательный аргумент, и~самопроверка гоняет оба.
\item \emph{Нулевая ошибка по~построению}: столбец «в~долях собственного здорового уровня» в~здоровой ячейке равен единице у~каждого набора, разброс нулевой, и~любой вердикт объявлялся достоверным.
\end{enumerate}

\textbf{Воспроизводимость.} Счёт разнесён по~файлам: модель, статистики, правила, сборщик блоков. Каждое правило \emph{объявляет}, какие блоки ему нужны; одинаковые блоки разных правил склеиваются по~хешу и~считаются один раз (из~78 запрошенных считается 42). Результат блока кладётся в~кэш, в~ключ которого входит хеш исходного кода модели: правка модели сама обесценивает кэш, и~старые блоки не~подмешаются к~новым. Отдельный тест требует \emph{побитового} совпадения с~прежним монолитом на~закреплённых точках. Вердикты пишутся в~журнал с~датой и~хешами блоков и~после счёта не~пересматриваются: негодность правила отмечается рядом, а~не~исправляется задним числом.

\begin{ineq}
Полезная мелочь, которая здесь дорого стоила. Размах $n$~значений величины с~относительной ошибкой~$r$ сам по~себе есть случайная величина: для~трёх независимых значений $\mathbb{E}[\text{размах}]=1{,}693\,r$. Поэтому порог «разброс по~$J$ не~больше 10~\%» осмыслен только при~$r<0{,}059$. При~300~наборах $r=12$~\%, то~есть правило было непроверяемо; при~2000 оно стало проверяемым~--- и~тут~же не~прошло, обнаружив третий путь, которым $J$ входит в~модель.
\end{ineq}
